\scanpage{644}
% 644-01
$\unsure{v|u}$ in $\mathrm{Gl}(2,\widehat{\mathbb Z})/\operatorname{Im}\widehat{\Pi}_0^{\unsure{\circ}}$, \unsure{factored} through $\mathcal G/\Pi_0$, is equal to that of $u$. Thus the assertion that $M^\circ(-\bar\jmath)$ is in the kernel of $M\longrightarrow\mathcal G/\Pi_0$, is equivalent \deleted{to the fact}\inserted{to the assertion} that $M^\circ(0)$ is in this kernel, hence that we have a commutative diagram
% 644-02
\diagram{\begin{tikzcd}[ampersand replacement=\&]
M(0) \arrow[r] \arrow[d,"{(\varepsilon_2,\varepsilon_3)}"'] \& M \arrow[d,"{(45)}"] \\
\mu\times\mu \arrow[r,"?"] \& \mathcal G/\Pi_0
\end{tikzcd}}
% 644-03
\margin{\deleted{\illegible{2}} N.B. \deleted{\unsure{commutative}; and} the \unsure{correction|condition} $\varepsilon_3$ does not enter into the def\unsure{inition} (47) of $M\longrightarrow\mathcal G/\Pi_0$, cf. further on.}
% 644-04
Note that the composite
% 644-05
\diagram{\begin{tikzcd}[ampersand replacement=\&, column sep=large,nodes in empty cells]
M \arrow[r] \arrow[dr] \& \mathcal G/\Pi_0 \arrow[r] \& (\deleted{\unsure{\mathrm{GL}}}\mathbb Z/4\mathbb Z)^*\simeq\{\pm1\} \\
\& \mathrm{Gl}(2,\mathbb Z/4\mathbb Z)^{\unsure{*}} \arrow[ur] \& \deleted{\unsure{\mathrm{GL}(2,\mathbb Z),\mathbb Z}} \arrow[l,shift left=1ex] \arrow[u,shift left=1ex]
\end{tikzcd}}
% 644-06
is none other than $\varepsilon_2$. Note also that by the compatibility
% 644-07
\diagram{\begin{tikzcd}[ampersand replacement=\&]
\mathcal G \arrow[r,hook] \arrow[dr,"{\mathrm{can.}}"'] \& M \arrow[d,"{(45)}"] \\
\& \mathcal G/\Pi_0
\end{tikzcd}}
% 644-08
the operation of $M$ on $\mathcal G/\Pi_0$ deduced from the homomorphism (45) is none other than the operation \deleted{\unsure{regular}}\inserted{\unsure{of a topological group|of quotients}}, of \unsure{automorphisms} on $\mathcal G$ on passing to the quotient by $\Pi_0$, an operation which, for $u\in M^!$ and in particular on $M(\unsure{\omega|0})\subset M^!$, is given by $\varepsilon_2(u)$ (cf. page \unsure{594|554}). Thus
% 644-09
\mspage{643}

\scanpage{645}
% 645-01
\mspage{625}
% 645-02
We see that $M_{\deleted{\unsure{0|a}}}^{\prime\prime!}$ maps into the center of $\unsure{\mathcal G/\Pi_0|\mathcal G[\Pi_0]}$, which is contained in $\mathcal G^+/\Pi_0$, and in fact equal to \deleted{center} the image $\mu$ of $\Pi/\Pi_0$ in $\mathcal G^+/\Pi_0$ (since the center of $\mathcal G^+/\Pi\simeq\mathcal G_3$ is trivial). Thus we have a factorization
% 645-03
\diagram{\begin{tikzcd}[ampersand replacement=\&]
M_{\deleted{\unsure{0|a}}}^{\prime\prime!} \arrow[r] \& \mu \arrow[r,hook,"{\mathrm{centre}}"'] \& \mathcal G/\Pi_0
\end{tikzcd}}
% 645-04
and we would have to verify that this homomorphism $M^{\prime\prime!}\longrightarrow\mu$ is none other than $\varepsilon_3$, which is therefore trivial on $M^{\circ!}$ and a fortiori on $M^{\circ!}(0)$.
% 645-05
\inserted{In fact, we \unsure{recognize} that the kernel is not $\operatorname{Ker}(\varepsilon_3\mid M^{\prime\prime!})$ \deleted{but} (which in \unsure{fact|fact} contains $\mu$;) but $M_0^{\prime\prime!}$ --- we will \uline{return} to this below.}
% 645-06
\deleted{Let us provisionally admit this \unsure{fact|point}, \unsure{find|let us find} then a factorization of
% 645-07
\diagram{\begin{tikzcd}[ampersand replacement=\&]
N_\rho^* \arrow[r] \& \deleted{\unsure{M}}\mathcal G/\Pi_0
\end{tikzcd}}
as
% 645-08
\diagram{\begin{tikzcd}[ampersand replacement=\&]
(*)\quad N_\rho^*/M^\circ(-\bar\jmath) \arrow[r] \& \mathcal G/\Pi_0
\end{tikzcd}}
and since the composite homomorphism
% 645-09
\diagram{\begin{tikzcd}[ampersand replacement=\&]
\mathbb D_\rho \arrow[r,hook] \& N_\rho^*/M^\circ(-\bar\jmath) \arrow[r] \& \mathcal G/\Pi_0
\end{tikzcd}}
% 645-10
is injective, we see that the kernel of (*) is a subgroup of $N_\rho^*/M^\circ(-\bar\jmath)$ whose intersection with $\mathbb D_\rho$ is $\{1\}$ --- it is therefore $1$, or of order $2$ (since}
% 645-11
\mspage{644}

\scanpage{646}
% 646-01
\deleted{$\mathbb D_\rho$ has index $2$ in $N_\rho^*/M^\circ(-\bar\jmath)$).
% 646-02
I see more clearly --- let us try the
% 646-03
\uline{Proposition} Consider the filtration of $\mathcal G^{\unsure{\wedge|*}}/\Pi_0^{\unsure{\wedge}}$ by $\mathcal G/\mathcal G^+\simeq\mu$, $\mathcal G^+/\Pi\simeq\mathcal G_3$, and $\Pi/\Pi_0\simeq\mu$. Consider the filtration of $M$ by $M/M^{\prime\prime}\xrightarrow[\sim]{\varepsilon_2}\mu$, by $M^{\prime\prime}/M^{\prime\prime!}\simeq\mathcal G_3$, and $M^{\prime\prime!}/M^{\circ!}\overset{\varepsilon_3}{\simeq}\mu$.
% 646-04
Then the canonical homomorphism
\diagram{\begin{tikzcd}[ampersand replacement=\&]
M \arrow[r] \& \mathcal G/\Pi_0
\end{tikzcd}}
% 646-05
is compatible with the filtrations, and induces \unsure{some} obvious isomorphisms on the factors. In particular the kernel of $M\longrightarrow\mathcal G/\Pi_0$ is $M^{\circ!}$.}
% 646-06
Admitting this point, we therefore find
% 646-07
\[
\operatorname{Ker}(N_\rho^*\longrightarrow\mathcal G/\Pi_0)
= N_\rho^*\cap M_{\inserted{0}}^{\deleted{\unsure{\circ}}\inserted{\prime\prime}!}\deleted{\unsure{a}}
\]
% 646-08
and since
\[
N_\rho^*\cap M_0^! = Z_{\rho0}^! = M^{\prime}(-\bar\jmath),
\]
% 646-09
we find that the kernel sought is equal to:
\[
M^{\prime}(-\bar\jmath)\cap M^{\prime\prime}=M^\circ(-\bar\jmath),\quad\mathrm{cqfd}.
\]
% 646-10
Thus \deleted{the homomorphism} we do have an injective homomorphism (46), and since the groups
% 646-11
\mspage{645}

\scanpage{647}
% 647-01
\mspage{626}
% 647-02
in question have the same order $24$, this homomorphism is indeed an isomorphism, q.e.d.
% 647-03
d) Application to the structure of $\mathcal M_{\rho,\sigma}$ and to the definition of a homomorphism \deleted{$\unsure{\mathcal M}_{\rho,\sigma}$} $M\longrightarrow\mathcal M_{\rho,\sigma}/\mu_\rho$.
% 647-04
Recall that
\[
\tag{48}
\mathcal M_{\rho,\sigma}\overset{\mathrm{d\acute ef}}{=}
N_\rho^*\times_{\gamma_{\rho,\sigma}}N_\sigma^*\times_{\gamma_{\rho,\sigma}}G_{\rho,\sigma}
\]
% 647-05
where \deleted{$\unsure{\mathcal G}_{\rho,\sigma}$}
\[
\tag{49}
\left\{\begin{aligned}
G_{\rho,\sigma}&\simeq M=\underbrace{M(0)}_{\inserted{\simeq M_{0,3}^{\sim}(0)}}\cdot\mathcal G^{+\wedge}\\
\gamma_{\rho,\sigma}&\overset{\mathrm{dfn}}{=}G_{\rho,\sigma}/\underbrace{\mathcal G^{+\wedge}}_{\inserted{\mathcal G_{\rho,\sigma}}}\simeq\widetilde{\Gamma}_{\mathbb Q}\ (\xleftarrow{\sim}M(0)),\\
N_\rho^*&=\operatorname{Norm}_M(L_\rho),\qquad N_\sigma^{*}=\operatorname{Norm}_M(L_\sigma)
\end{aligned}\right.
\]
% 647-06
On the other hand, we have introduced the distinguished subgroups
\[
\tag{50}
\left\{\begin{aligned}
Z_{\sigma0}^!&=M(\tfrac12)\subset N_\sigma^*\qquad(\mathrm{p.}\ 615\ (9))\\
\deleted{\unsure{A}}\inserted{N}_\rho^?&=\deleted{\unsure{\rho}}Z_\rho^!\cdot\mu_{\tau'}\subset N_\sigma^*\qquad(\mathrm{p.}\ 616\ (21)).
\end{aligned}\right.
\]
% 647-07
The first of these groups does not contain $\mu$, the second contains it
\[
\tag{51}
\left\{\begin{aligned}
Z_{\sigma0}^!\cap\mu&=\{1\}\\
N_\rho^?\deleted{\unsure{0}}&\supset\mu
\end{aligned}\right.
\]
% 647-08
and
\diagram{(52)\quad$\left\{\begin{array}{l}
\begin{tikzcd}[ampersand replacement=\&]
Z_{\sigma0}^! \arrow[r,"\sim"] \& \gamma_{\rho,\sigma}=\widetilde{\Gamma}_{\mathbb Q}
\end{tikzcd}\\[1.2em]
\deleted{\unsure{\operatorname{Ker}}\,N_\rho^?\unsure{=\mu|\longrightarrow\mu}}\quad
\begin{tikzcd}[ampersand replacement=\&, column sep=small]
{} \& {} \& {} \& \inserted{\widetilde{\Gamma}_{\mathbb Q}}\arrow[d,equal] \& \text{exact}\arrow[d,bend left=25]\\
1\arrow[r] \& \mu\arrow[r] \& N_\rho^?\arrow[r] \& \gamma_{\rho,\sigma}\arrow[r] \& 1
\end{tikzcd}
\end{array}\right.$}
% 647-09
\mspage{646}

\scanpage{648}

Now let
\[
\begin{array}{rcl}
(53)\qquad \mathcal{M}_{\rho,\sigma}^{\unsure{\natural|\sharp}}
&\subset&\mathcal{M}_{\rho,\sigma}\\[-2pt]
\Vert\,\mathrm{dfn}&&\\[-2pt]
N_{\rho}^{?}\times_{\gamma_{\rho,\sigma}}
Z_{\sigma0}^{!}\times_{\gamma_{\rho,\sigma}}
\underbrace{G_{\rho,\sigma}}_{M}&&
\end{array}
\]

We have a canonical \deleted{\unsure{homomorphism}}\inserted{central embedding}
\[
\deleted{\epsilon\mapsto(1,\epsilon,1)}
\]
\diagram{
\begin{tikzcd}[ampersand replacement=\&, column sep=large]
(54)\qquad \mu\times\mu\times\mu \arrow[r,hook] \& \mathcal{M}_{\rho,\sigma}\\
(\epsilon,\epsilon',\epsilon'') \arrow[r,mapsto] \& (\epsilon,\epsilon',\epsilon'')
\end{tikzcd}}
coming from the central embeddings
\deleted{\unsure{the}}
\diagram{
\begin{tikzcd}[ampersand replacement=\&, column sep=small]
\mu \arrow[r,hook] \& N_{\rho}^{*}, \& \mu \arrow[r,hook] \& N_{\sigma}^{*}, \&
\mu \arrow[r,hook] \arrow[dr,hook,bend right=25] \& G_{\rho,\sigma}\\
\&\&\&\&\& \mathcal{G}_{\rho,\sigma}^{+}=\mathcal{G}^{+}\unsure{)} \arrow[u,hook,bend right=12]
\end{tikzcd}}
and we have
\[
\mathcal{M}_{\rho,\sigma}^{\unsure{\natural|\sharp}}
\cap(\mu\times\mu\times\mu)=(1\times\mu\times\mu)
\]
hence
\diagram{
\begin{tikzcd}[ampersand replacement=\&]
(55)\qquad \mu\times\mu \arrow[r,hook] \& \mathcal{M}_{\rho,\sigma}^{\unsure{\natural|\sharp}}.
\end{tikzcd}}

We denote by $\mu_{\rho}$ the central subgroup $\mu$ of $N_{\rho}$ or of $N_{\rho}^{?}$,
or $\mu\times1$ in $\mathcal{M}_{\rho,\sigma}^{\unsure{\natural|\sharp}}$.
It is then clear from (52), that we have an exact sequence
\diagram{
\begin{tikzcd}[ampersand replacement=\&, column sep=large]
(55)\qquad 1 \arrow[r] \& \mu_{\rho} \arrow[r,hook] \&
\mathcal{M}_{\rho,\sigma}^{\unsure{\natural|\sharp}} \arrow[r] \&
G_{\rho,\sigma} \arrow[r] \arrow[d,phantom,"\simeq"{rotate=90}] \& 1\\
\&\&\&M\&
\end{tikzcd}}
whence canonical \unsure{homomorphisms|homomorphisms}

\mspage{647}

\scanpage{649}
\mspage{627}

\diagram{
\begin{tikzcd}[ampersand replacement=\&, column sep=large]
(56)\qquad M \arrow[r,"\sim"] \&
\mathcal{M}_{\rho,\sigma}^{\unsure{\natural|\sharp}}/\mu_{\rho} \arrow[r,hook] \&
\mathcal{M}_{\rho,\sigma}/\mu_{\rho}.
\end{tikzcd}}

\margin{We therefore set
$\Gamma_{\rho,\sigma}^{\unsure{\natural|\sharp}}
\overset{\mathrm{dfn}}{=}
\mathcal{M}_{\rho,\sigma}^{\unsure{\natural|\sharp}}/\mathcal{G}_{\rho,\sigma}^{+}
\subset\Gamma_{\rho,\sigma}=\mathcal{M}_{\rho,\sigma}/\mathcal{G}_{\rho,\sigma}^{+}$.}

Obviously the subgroup
$\mathcal{M}_{\rho,\sigma}^{\unsure{\natural|\sharp}}$ of $\mathcal{M}_{\rho,\sigma}$
contains $\mathcal{G}_{\rho,\sigma}=SG_{\rho,\sigma}\unsure{\simeq}\mathcal{G}^{+\wedge}$,
in fact we have a diagram of exact sequences
\diagram{
\begin{tikzcd}[ampersand replacement=\&, column sep=small,row sep=large]
1 \arrow[r] \& \mu_{\rho}\times\mathcal{G}_{\rho,\sigma}^{\deleted{\unsure{+}}}
\arrow[r] \arrow[d,hook] \&
\mathcal{M}_{\rho,\sigma}^{\unsure{\natural|\sharp}} \arrow[r] \arrow[d] \&
\gamma_{\rho,\sigma} \arrow[r] \arrow[d,equal] \&1\\
(57)\quad1 \arrow[r] \&
\underbrace{SZ_{\rho}}_{L_{\rho}}\times\underbrace{SZ_{\sigma}}_{L_{\sigma}}
\times\mathcal{G}_{\rho,\sigma}^{+} \arrow[r] \&
\mathcal{M}_{\rho,\sigma} \arrow[r] \& \unsure{\gamma_{\rho,\sigma}|\Gamma_{\rho,\sigma}} \arrow[r] \&1
\end{tikzcd}}

In the isomorphism (56), the subgroup $\mathcal{G}^{+\wedge}$ of $M$
corresponds to the subgroup
$\mu_{\rho}\times\mathcal{G}_{\rho,\sigma}^{+}/\mu_{\unsure{\sigma|\rho}}
\simeq\mathcal{G}_{\rho,\unsure{\sigma}}$,
it maps isomorphically onto the subgroup $\mathcal{G}_{\rho,\sigma}$
of $\mathcal{M}_{\rho,\sigma}/\mu_{\rho}$.

\deleted{The subgroup $M(0)$ maps}

We thus find
\diagram{
\begin{tikzcd}[ampersand replacement=\&, column sep=small,row sep=large]
\& \inserted{\unsure{\mathcal{G}_{\rho,\sigma}^{+}|G_{\rho,\sigma}^{+}}}
\arrow[d,phantom,"\simeq"{rotate=90}] \&\&
\inserted{\gamma_{\rho,\sigma}} \arrow[d,phantom,"\simeq"{rotate=90}]\&\\
1 \arrow[r] \& \mathcal{G}^{+\wedge} \arrow[r] \arrow[d,"\sim"] \&
M \arrow[r] \arrow[d,"\sim"] \& \widetilde{\Gamma}_{\mathbb{Q}}
\arrow[r] \arrow[d,phantom,"\simeq"{rotate=90}] \&1\\
(58)\quad1 \arrow[r] \& \mathcal{G}_{\rho,\sigma}^{+}
\arrow[r,"\unsure{\sim}"] \arrow[d,hook,"1"] \&
\mathcal{M}_{\rho,\sigma}^{\unsure{\natural|\sharp}}/\mu_{\rho} \arrow[r] \arrow[d,hook] \&
\Gamma_{\rho,\sigma}^{\unsure{\natural|\sharp}}/\mu_{\rho} \arrow[r] \arrow[d,hook] \&1\\
1 \arrow[r] \& \mathcal{G}_{\rho,\sigma}^{+} \arrow[r] \&
\mathcal{M}_{\rho,\sigma}/\mu_{\rho} \arrow[r] \& \Gamma_{\rho,\sigma}/\mu_{\rho} \arrow[r] \&1
\end{tikzcd}}

Thus, we find in passing that the canonical \unsure{homomorphism}
\diagram{
\begin{tikzcd}[ampersand replacement=\&, column sep=small,row sep=large,nodes in empty cells]
(59)\quad1 \arrow[r] \&
\underbrace{SZ_{\rho}\times SZ_{\sigma}}_{L_{\rho}\times L_{\sigma}} \arrow[r] \&
\Gamma_{\rho,\sigma} \arrow[r,"\inserted{\mathrm{\acute{e}pi}}"{yshift=3pt}] \&
\gamma_{\rho,\sigma} \arrow[r] \&1\\
\& \text{admits a ``}\unsure{\text{bisection}}\text{''} \&\&
\deleted{\unsure{\Gamma_{\rho,\sigma}^{\natural}/\mu_{\rho}}} \arrow[u,"\sim"] \&
\Gamma_{\rho,\sigma}^{\unsure{\natural|\sharp}}\text{ (whose}
\end{tikzcd}}
intersection with $\unsure{SZ_{\rho}}\times SZ_{\sigma}$ is $\mu_{\rho}\times1$)

\mspage{648}

\scanpage{650}

\[
(60)\qquad
\left\{\begin{aligned}
\Gamma_{\rho,\sigma}&\simeq
\Gamma_{\rho,\sigma}^{\unsure{\natural|\sharp}}\cdot
\underbrace{S\Gamma_{\rho,\sigma}}_{SZ_{\rho}\times SZ_{\sigma}=L_{\rho}\times L_{\sigma}}
\ ;\quad
\Gamma_{\rho,\sigma}^{\unsure{\natural|\sharp}}\cap S\Gamma_{\rho,\sigma}=\mu_{\rho}
\\
\Gamma_{\rho,\sigma}/\mu_{\rho}&\simeq
\left(\underbrace{\Gamma_{\rho,\sigma}^{\unsure{\natural|\sharp}}/\mu_{\rho}}
_{\gamma_{\rho,\sigma}\simeq\widetilde{\Gamma}_{\mathbb{Q}}}\right)
\cdot\left(L_{\rho}/\mu_{\rho}\times L_{\sigma}\right).
\end{aligned}\right.
\]

To say what becomes of the remarkable subgroups
\[
M(0),\quad
\begin{array}{c}
M[0],\\
\inserted{\Vert}\\[-4pt]
\inserted{M^{\mathcal{G}}[0]}\\
\inserted{\Vert}\\[-4pt]
\inserted{M(\unsure{\omega|0}).L_{0}^{\mathcal{G}}}
\end{array}
\quad
\begin{array}{c}
M_{\unsure{a}}^{!}[0],\\
\inserted{\Vert}\\[-4pt]
\inserted{M^{\unsure{\Pi}}[0]}\\
\inserted{\Vert}\\[-4pt]
\inserted{M(\unsure{\omega|0}).L_{0}^{\Pi}}
\end{array}
\quad
\begin{array}{c}
M_{0}^{\unsure{!}}[0]\\
\inserted{\Vert}\\[-4pt]
\inserted{M^{\Pi_{0}}\unsure{[0]}}\\
\inserted{\Vert}\\[-4pt]
\deleted{\unsure{M(\omega)L_{0}|M(0)L_{0}}}\\
\inserted{M(\unsure{\omega|0})L_{0}^{\Pi_{0}}}
\end{array}
\]
\[
\begin{array}{c}
\unsure{\{M(\tfrac12)\}|M(\tfrac12)}\\
\inserted{\Vert}\\[-4pt]
\inserted{Z_{\sigma0}^{!}}
\end{array},\quad
M'(-\bar{\jmath}),\quad M'''(-\bar{\jmath}),\quad N_{\rho}^{?}
\]
of $M$ under the homomorphism
\diagram{
\begin{tikzcd}[ampersand replacement=\&, column sep=large]
M \arrow[r,"\sim"] \&
\mathcal{M}_{\rho,\sigma}^{\unsure{\natural|\sharp}}/\mu_{\rho} \arrow[r,hook] \&
\mathcal{M}_{\rho,\sigma},
\end{tikzcd}}
we would first have to introduce the appropriate groups in the general $\mathcal{M}_{\rho,\sigma}$.
We are therefore going to return to these groups now.

\mspage{649}

\scanpage{651}
\mspage{628}

\uline{VI}\quad The general groups $\mathcal{M}_{\rho,\sigma}$:
recapitulation, refinement.

$1^\circ)$ Data. These are those of \uline{V} $1^\circ)$ (page~577).
In addition to $\mathcal{G}$ \inserted{(profinite)} and to
\diagram{\begin{tikzcd}[ampersand replacement=\&]
\mathrm{Gl}(2,\mathbb{Z})^{\wedge} \arrow[r] \& \mathcal{G}
\end{tikzcd}}
a continuous surjective homomorphism, i.e. generators $\rho,\sigma,\tau$
(or $\rho,\sigma,\tau'$) of $\mathcal{G}$, we also give ourselves an open subgroup
\[
(1)\qquad \mathcal{G}_{\unsure{0|a}}^{\unsure{\natural|\sharp}}
\subset\mathcal{G}^{+},\quad
\text{invariant in }\mathcal{G}^{\unsure{*|+}}.
\]

In the universal case $\mathcal{G}=\mathrm{Gl}(2,\mathbb{Z})^{\wedge}$,
the datum of $\mathcal{G}_{\unsure{0|a}}^{\unsure{\natural|\sharp}}$
may correspond to the point of view where we look at the subgroups
$M^{\unsure{\natural|\sharp}}$ of $M$
\deleted{which}\inserted{\unsure{formed|formed} of the $u$ which}
\unsure{normalize}\deleted{\unsure{ateur of|ent the}}
$\mathcal{G}^{\unsure{\natural|\sharp}}$
\inserted{and which operate \unsure{trivially} in
$\mathcal{G}^{+}/\mathcal{G}^{\unsure{\natural|\sharp}}$ --- this is}
\deleted{--- which is \unsure{necessarily} open, and defines}
\inserted{an open subgroup, which defines}
in $\widetilde{\Gamma}_{\mathbb{Q}}=M/\mathcal{G}^{+}$
an open image subgroup
$\widetilde{\Gamma}_{\mathbb{Q}}^{\unsure{\natural|\sharp}}$,
whose representations in suitable profinite groups we may propose to study.
An interesting case will be the one where
\[
\mathcal{G}^{\unsure{\natural|\sharp}}=\Pi_{0},\quad
\text{thus }M^{\unsure{\natural|\sharp}}=M^{\prime\prime!},
\]
or
\[
\mathcal{G}^{\unsure{\natural|\sharp}}=\Pi,\quad
\text{thus }M^{\unsure{\natural|\sharp}}=M^{!}.
\]
If
\[
\mathcal{G}^{\unsure{\natural|\sharp}}=\mathcal{G},\quad
\text{thus }M^{\unsure{\natural|\sharp}}=M,
\]
we recover the preceding constructions.
\deleted{\unsure{We can}}

\mspage{650}

\scanpage{652}

An interesting case also is the one where $\unsure{\mathcal{G}}$ is the group
of adelic \inserted{integral} points of a group scheme over $\mathbb{Z}$
(in this case, it is true that $\rho$ and $\sigma,\tau$ generate
\unsure{perhaps|not then} $\mathcal{G}^{\unsure{+|*}}$ only modulo a finite group
i.e. they generate \unsure{perhaps} only an open subgroup.
In this case, we should rather \uline{\unsure{paraphrase}} the constructions
that follow in the scheme-theoretic case
I shall come back to this later on. An interesting choice of a
$\mathcal{G}^{\unsure{\natural|\sharp}}$ is then the
\uline{identity component} of $\mathcal{G}$.

I hesitate to introduce notations
$Z_{\rho,\sigma}^{\unsure{\natural|\sharp}},\mathcal{M}_{\rho,\sigma}^{\unsure{\natural|\sharp}}$
etc --- so as not to overload notations that are already heavy. So I prefer to
use the simpler notations $Z_{\rho,\sigma}$, etc --- it being understood
that in these constructions, the group $\mathcal{G}^{\unsure{\natural|\sharp}}$
enters --- with the option of reintroducing the notation $\unsure{\natural|\sharp}$
if need be.
We put
\[
(3)\qquad
\mathcal{G}_{\tau}^{\unsure{\natural|\sharp}}
=\overbrace{\{1,\tau\}}^{\mu_{\tau}}\,.
\mathcal{G}^{\unsure{\natural|\sharp}}\subset\mathcal{G}.
\]

\mspage{651}

\scanpage{653}
\mspage{629}

\Needspace{14\baselineskip}
$2^\circ)$ \quad $Z_{\rho,\sigma},\ \gamma_{\rho,\sigma}$.

Let
\[
\arraycolsep=3pt (4)\qquad
\begin{array}{l}
 Z_{\rho,\sigma}\subset
 \operatorname{Aut}(\mathcal{G}^{+})\times
 \widehat{\mathbb{Z}}^{*}\times\mu\times\mu\\
 \phantom{Z_{\rho,\sigma}}\Vert\,\mathrm{dfn}\\
 \left\{(u,\mu,\varepsilon_2,\varepsilon_3)\ \left|
 \begin{array}{ll}
 \multicolumn{2}{l}{\deleted{\unsure{u\ \text{normalizes}\ \mathcal{G}^{\natural}\ 
 \text{and induces the identity in}\ \mathcal{G}^{+}/\mathcal{G}^{\natural}}}}\\
 a)&\exists\alpha\in\mathcal{G}_{\tau}^{\unsure{\natural|\sharp}},
 \ \text{with }\varepsilon_{\mathcal{G}}(\alpha)=\varepsilon_3,
 \quad u(\rho)=\alpha(\rho)\\
 b)&\exists\beta\in\mathcal{G}_{\tau}^{\unsure{\natural|\sharp}},
 \ \text{with }\varepsilon_{\mathcal{G}}(\beta)=\varepsilon_2,
 \quad u(\sigma)=\beta(\sigma)\\
 c)&u(\varepsilon_0)=\varepsilon_0^\mu\\
 d)&u\ \text{normalizes }\mathcal{G}^{\unsure{\natural|\sharp}}
 \inserted{\ \text{\uline{and} induces the identity on }
 (\mathcal{G}^{+}/\mathcal{G}^{\unsure{\natural|\sharp}})/\operatorname{Im}\Pi}
 \end{array}\right.\right\}.
\end{array}
\]
\margin{\unsure{cf. page} \unsure{635} \unsure{continuation}
\deleted{\unsure{a weaker condition}}
\unsure{this condition}.}
\margin{\deleted{The last condition \unsure{added} in d)
(\unsure{in principle}), is a \unsure{consequence} of a)b) if
$\varepsilon_2\varepsilon_3=1$
(\unsure{in which case} $u|(\mathcal{G}^{+}/\mathcal{G}^{\unsure{\natural|\sharp}})$
\deleted{the identity} \unsure{induced by} $\varepsilon[\tau]$,
\unsure{where} $\varepsilon=\varepsilon_2=\varepsilon_3\in\{\pm1\}$)}.}

\uline{Proposition.} \quad $Z_{\rho,\sigma}$ is a \uline{subgroup}
of the product group.

\uline{Proof.} The conditions c), d) define \unsure{obviously}
\deleted{a} \inserted{each} subgroup.
\unsure{Given that}
\[
\tau(\sigma)=-\unsure{\sigma|\varepsilon},
\]
condition b) also means that
\[
\unsure{(5^{*})}\qquad
b')\quad u(\sigma)\ \text{is }\mathcal{G}^{\unsure{\natural|\sharp}}
\text{-conjugate to : }\
\deleted{\unsure{\varepsilon_2}}\sigma^{\varepsilon_2}
\bigl(\deleted{\unsure{\varepsilon_2\sigma}}\inserted{=\varepsilon_2\sigma}\bigr)
\]
and in this form it is clear that condition c')
defines a subgroup. I do not think that condition a)
by itself defines a subgroup. But if
$\unsure{\mathfrak{g}_{(b)}|G_{(b)}}$
\deleted{$\unsure{\mathfrak{g}_{(a,b)}}$}
denotes the sub\unsure{group} of the product defined by these conditions
\deleted{\unsure{a)}}\inserted{b)}, we have \unsure{to show} that
$\unsure{\mathfrak{g}|G}_{\deleted{\unsure{(a,b,c)}}\inserted{(a,b)}}$
is a subgroup of the \deleted{product}
$\unsure{\mathfrak{g}|G}_{\deleted{\unsure{(a,b)}}\inserted{(b)}}$
--- which will prove the proposition.
Let $\uline{u},\uline{u}'\in\unsure{\mathfrak{g}_{(b)}|G_{(b)}}$,
corresponding to
\inserted{$(u,\mu,\varepsilon_2,\varepsilon_3)$ and
$(u',\mu',\varepsilon'_2,\varepsilon'_3)$, and to}
$\alpha,\alpha'\in\mathcal{G}^{\unsure{\natural|\sharp}}_\tau$,
\unsure{so} that we have

\mspage{652}

\scanpage{654}

\[
(6)\qquad
u(\rho)=\alpha(\rho),\quad u'(\rho)=\alpha'(\rho),
\inserted{\quad\text{with }\begin{array}{l}
\varepsilon_{\mathcal{G}}(\alpha)=\varepsilon_3,\\
\alpha_{\mathcal{G}}(\alpha')=\varepsilon'_3.
\end{array}}
\]
\unsure{Let us prove} \deleted{$\unsure{\varepsilon_3}$} that
$u'u\in\unsure{\mathfrak{g}_{(a)(b)}|G_{(a,b)}}$
--- i.e. that it also satisfies \unsure{(b)|(a)}.
If $\varepsilon_3=1$ i.e.
$\alpha\in\mathcal{G}^{\unsure{\natural|\sharp}}
\inserted{\subset\mathcal{G}^{+}}$,
$u'(\alpha)$ is defined and we have
\[
\begin{aligned}
(u'u)(\rho)&=u'(u(\rho))=u'(\alpha(\rho))
=u'(\alpha)(u'(\rho))=u'(\alpha)(\alpha'(\rho))\\
&=(u'(\alpha)\alpha')(\rho)=\alpha''(\rho),
\end{aligned}
\]
with
\[
(7)\qquad
\begin{array}{l|l}
\alpha''=u'(\alpha)\alpha'\in
\mathcal{G}^{\unsure{\natural|\sharp}}_\tau&
\text{Case }\varepsilon_3=1\\
\varepsilon_{\mathcal{G}}(\alpha'')
=\varepsilon_{\mathcal{G}}(\alpha')
=\varepsilon'_3=\varepsilon_3\varepsilon'_3&
\end{array}
\]
(since $u'(\alpha)\in\mathcal{G}^{\unsure{\natural|\sharp}}$
i.e. $\varepsilon_{\mathcal{G}}(u'(\alpha))=1$, and $\varepsilon_3=1$.

In this case we \unsure{win}. If $\varepsilon_3=-1$, we write
\[
\alpha=\alpha_0\tau\qquad
\alpha_0\in\mathcal{G}^{\unsure{\natural|\sharp}},
\]
and we have
\[
u(\rho)=(\alpha_0\tau)(\rho)=\alpha_0(\tau(\rho))
=\alpha_0\bigl(\sigma^{\unsure{-1}}(\rho^{-1})\bigr)
=(\alpha_0\sigma)(\rho^{-1}),
\]
where now $\alpha_0\sigma\in\mathcal{G}^{+}$, hence
$u'(\alpha_0\sigma)$ is defined, and we will have
\[
\unsure{(u'u)(\rho)}
=u'(u(\rho))
=u'(\alpha_0\sigma)
\bigl(\underbrace{u'(\rho^{-1})}_{\alpha'(\rho^{-1})}\bigr)
=\bigl(u'(\alpha_0\sigma)\alpha'\bigr)(\rho^{-1})
\]
and since
\[
\deleted{\unsure{\rho^{-1}}}\qquad
\rho^{-1}=\sigma\tau(\rho)
\]
we find
\[
(8)\qquad
(u'u)(\rho)=
\deleted{\unsure{u'(\alpha_0\sigma)\alpha'\sigma\tau}}
\alpha''(\rho)
\]
with
\[
\left\{
\begin{array}{l}
\alpha''=\varepsilon
\overbrace{u'(\alpha_0)}^{\inserted{\in\mathcal{G}^{\unsure{\natural|\sharp}}}}
\overbrace{u'(\sigma)}^{\inserted{\in\mathcal{G}^{+}}}
\overbrace{\alpha'}^{\inserted{\in\mathcal{G}^{\unsure{\natural|\sharp}}_\tau}}
\underset{\deleted{\unsure{\tau'}}}{\sigma\tau}
\qquad\text{where}\\
\varepsilon\in\mu=\{\pm1\}
\unsure{\hookrightarrow|\Rightarrow}\operatorname{Centre}\mathcal{G}^{+}
\end{array}
\right.
\]
\deleted{\unsure{is}} is chosen so that $\alpha''$ (if possible)
is in $\mathcal{G}^{\unsure{\natural|\sharp}}_\tau$.
Let us note that all the factors of this product are in
$\unsure{\mathcal{G}^{\natural}_\tau\mathcal{G}^{+}|
\mathcal{G}^{\natural}_\tau/\mathcal{G}^{+}}$,
\unsure{\uline{in}} any case $\alpha''\in\mathcal{G}$.

\mspage{653}

\scanpage{655}
\mspage{630}

Let us reduce modulo $\mathcal{G}^{\unsure{\natural|\sharp}}$,
the relation $\alpha''\in\mathcal{G}^{\unsure{\natural|\sharp}}_\tau$
reads
\[
\dot{\alpha}''\in\{1,\dot{\tau}\}\quad\text{i.e.}\quad
\dot{\varepsilon}\,u'(\sigma)^{\cdot}\dot{\alpha}'\dot{\sigma}\dot{\tau}
\in\{1,\dot{\tau}\}
\quad\text{i.e.}\quad
\dot{\varepsilon}\,u'(\sigma)^{\cdot}\dot{\alpha}'\dot{\sigma}
\in\{1,\dot{\tau}\}.
\]
\deleted{\unsure{we would have}, since $\sigma\tau=-\tau\sigma$}
\inserted{We distinguish \unsure{again} two cases}

$1^\circ)$ \quad $\varepsilon'_3=1$ i.e. $\dot{\alpha}'=1$
\deleted{$\unsure{\dot{\varepsilon}\,u'(\sigma)^{\cdot}
\dot{\alpha}'\dot{\tau}}$}
the relation reads
\[
(*)\qquad \dot{\varepsilon}\,u'(\sigma)^{\cdot}\dot{\sigma}=1
\]
$2^\circ)$ \quad $\varepsilon'_3=-1$ i.e. $\dot{\alpha}'=\dot{\tau}$,
since $\tau\sigma=-\sigma\tau$, the relation reads
\[
(**)\qquad -\dot{\varepsilon}\,u'(\sigma)^{\cdot}\dot{\sigma}=1,
\]
thus in all cases, the relation reads
\[
(9)\qquad
\left\{
\begin{array}{l}
u'(\sigma)^{\cdot}=(\varepsilon\varepsilon'_3)^{\cdot}
\dot{\sigma}^{-1}
\quad\deleted{\unsure{\text{or}}}\ \text{i.e. again}\\
u'(\sigma)^{\cdot}=(-\varepsilon\varepsilon'_3)^{\cdot}\dot{\sigma}
\end{array}
\right.
\]
Now \deleted{let us use hypothesis b) on $u'$,}
\inserted{suppose now that we have b) for $u'$,}
which \deleted{\unsure{means}}\inserted{implies}
\[
u'(\sigma)^{\cdot}
=\varepsilon'_2\,\deleted{\text{\illegible{1}}}\dot{\sigma},
\]
we see therefore that it suffices to choose $\varepsilon$ so that
$-\varepsilon\varepsilon'_3=\varepsilon'_2$
i.e.
\[
\varepsilon=-\varepsilon'_2\varepsilon'_3
\]
thus we will have
\[
\deleted{\unsure{(10)}}\qquad
\alpha''=-\varepsilon'_2\varepsilon'_3\,
u'\bigl(
\underbrace{\alpha_0\deleted{\unsure{\tau}}\sigma}
_{\inserted{\alpha\tau\sigma=\alpha\tau'}}
\bigr)\alpha'\underbrace{\sigma\tau}_{\inserted{-\tau'}}
\]
i.e.
\[
(10)\qquad
\alpha''=\varepsilon'_2\varepsilon'_3\,u'(\alpha\tau')
(\alpha'\tau')
\quad\big|\quad\text{Case }\varepsilon_3=-1
\]
a formula which makes sense, since $\alpha\tau'\in\mathcal{G}^{+}$ hence
$u'(\alpha\tau')$ is defined, and the preceding argument proves that
the element in question \deleted{\unsure{is in}}\inserted{satisfies}
\[
(11)\qquad
\alpha''\in\mathcal{G}^{\unsure{\natural|\sharp}}_\tau,\quad
\varepsilon_{\mathcal{G}}(\alpha'')
=\underbrace{\varepsilon_{\mathcal{G}}(\alpha)}_{-1}
\varepsilon_{\mathcal{G}}(\alpha')
\]
by distinguishing cases \unsure{(*)} and $(**)$ of (9).
Same conclusion in the first case where
$\varepsilon_{\mathcal{G}}(\alpha)=1$ (7).
\inserted{This completes the proof.}

\mspage{654}

\scanpage{656}

\Needspace{10\baselineskip}
\uline{Remarks} \deleted{a) For condition b) to be
\unsure{multiplicative}, we do not need condition a) in
\unsure{all} its strength, it suffices that $u$ in
$\mathcal{G}^{+}/\mathcal{G}^{\unsure{\natural|\sharp}}$
normalize the image of $L_\sigma$, i.e. that we have
\[
u(\sigma)\equiv\pm\sigma
\quad(\mathcal{G}^{\unsure{\natural|\sharp}}).
\]}

\deleted{b)} If $\mathcal{G}$ is ``big enough'', the projection homomorphism
is injective

\begin{samepage}
(12)\quad
\diagram{\begin{tikzcd}[ampersand replacement=\&, column sep=large]
Z_{\rho,\sigma}\arrow[r,hook]\&\operatorname{Aut}(\mathcal{G}^{+})
\end{tikzcd}}
\end{samepage}
\quad (cf. page 578)

\uline{Examples} in the universal case

a) \quad $\mathcal{G}^{\unsure{\natural|\sharp}}=\mathcal{G}^{+}
=\deleted{\unsure{\mathcal{G}}}\mathrm{Sl}(2,\mathbb{Z})^{\wedge}$.
\deleted{Here we have the} \unsure{Condition|condition} (4) d) is empty,
conditions a), b) are equivalent respectively to the fact that
$u(\rho)$ conjugate to $\rho^{\varepsilon_3}$,
$u(\sigma)$ conjugate to $\sigma^{\varepsilon_2}$,
i.e. \unsure{comes back} to definition (6) page 578, which gives
\[
\deleted{Z_{\rho,\sigma}=M(\unsure{\omega|0}).L_0^{\mathcal{G}}
=M^{\mathcal{G}}[0].}
\]
\[
(13)\qquad
Z_{\rho,\sigma}=M(0).L_0^{\mathcal{G}}
=M^{\mathcal{G}}[0]=M[0]
\]

\margin{\uline{N.B.} A priori we rather have
$Z_{\rho,\sigma}\subset\unsure{\widetilde{M}_{0,3}|M_{0,3}}$, but
\diagram{\begin{tikzcd}[ampersand replacement=\&, column sep=large,row sep=large]
M[0]\arrow[r,"\sim"]\arrow[d,equal]\&
\widetilde{M}_{0,3}[0]\arrow[d,equal]\\
M(\unsure{\omega|0}).L_0^{\mathcal{G}}\arrow[r,"\sim"]\&
\unsure{M_{0,3}(0)}.
L_0^{\tikz[baseline=(g.base)]{\node[draw,ellipse,inner sep=1pt]
(g){$\scriptstyle\unsure{\mathcal{G}_{0,3}}$};}}
\end{tikzcd}}}

b) \quad
$\mathcal{G}^{\unsure{\natural|\sharp}}=\Pi=\Pi_0\times\mu$.
We now find
\[
(14)\qquad
Z_{\rho,\sigma}=M(0).L_0^\Pi
=M^\Pi[0]=\unsure{M^{!}[0]|M^{\prime!}[0]}
\]
(denoted $Z_{\rho,\sigma}^{!}$ in $\mathrm{V}$ number$^\circ$ 8).

\margin{In this case, it is immediate that $Z_{\rho,\sigma}$
operates \unsure{trivial\textsuperscript{ly}} even on
$\mathcal{G}^{+}/\mathcal{G}^{\unsure{\natural|\sharp}}_\tau$,
for this is so for the
\deleted{\unsure{N.B.} \unsure{all the subgroups}
\unsure{contained by} a), b) of (4),
$\unsure{\mathcal{G}^{\natural}=M(0)_2.L_0^\Pi|
\mathcal{G}^{\natural}=M(\omega)_2.L_0^\Pi}$}.}
\margin{$\unsure{M^{!}[0]|M^{\prime!}[0]}=M^\Pi[0]$.}
\margin{make explicit again why
\unsure{(i.e. why)} that
$(\varepsilon_0,1,1,1)\notin Z_{\rho,\sigma}$}

c) \quad $\mathcal{G}^{\unsure{\natural|\sharp}}=\Pi_0$.
Conditions a) and b) \inserted{of (4)} have not changed
\unsure{with respect to|with respect to the}
\deleted{\unsure{before}}\inserted{\uline{preceding case}},
since $\Pi=\Pi_0\times\mu$ and $\mu$ is central, on the other hand the

\mspage{655}

\scanpage{657}
\mspage{631}

condition d) becomes distinctly stricter, we find for example that
$Z_{\rho,\sigma}\cap\mathcal{G}^{+}=L_0^{\Pi_0}$
(of index 2 in the corresponding group $L_0^\Pi$ of the preceding case),
and since the operation of $M(0)$ \inserted{$\subset M^!$} on
$\mathcal{G}^{+}/\Pi_0$ is through the character
$\varepsilon_2^{\unsure{*|\#}}$, we see that we now have
\[
(15)\qquad
Z_{\rho,\sigma}=M''(0).L_0^{\Pi_0}=M''[0]_0
\]
\[
\left\{
\begin{array}{l}
M''(0)=\operatorname{Ker}(M(0)\xrightarrow{\varepsilon_2}\mu)\\
M''[0]_0=\operatorname{Ker}(M[0]_0\xrightarrow{\varepsilon_2}\mu)\\
M[0]_0=M(\unsure{\omega|0}).L_0^{\Pi_0}
\end{array}
\right.
\]

\margin{\deleted{This is so if in definition (4) of
$Z_{\rho,\sigma}$ we require in d) that $u$
\unsure{induce|produce} the trivial operation on
$\mathcal{G}^{+}/\mathcal{G}^{\unsure{\natural|\sharp}}$,
a condition that I ended up abandoning as useless.
\unsure{In the new version} the cases
$\mathcal{G}^{\unsure{\natural|\sharp}}=\Pi$ and
$\mathcal{G}^{\unsure{\natural|\sharp}}=\Pi_0$
give the same $Z_{\rho,\sigma}$ (15).
But I would like to obtain
\[
M^{\Pi_0}[0]=\unsure{M[0]^!_0|M[0]_0}
=M(\unsure{\omega|0}).\Pi_0
\]
--- the group of b) is \unsure{too} large, that of c)
\unsure{a little too}
\uline{\unsure{(old version)}}
\inserted{a little too small!}}}

\margin{\deleted{But with the new definition of
$Z_{\rho,\sigma}$ (condition added to (4)) we will finally have
\[
Z_{\rho,\sigma}=M^![0]
\]
as in the preceding case b).
\inserted{We could have \unsure{introduced}
an $\unsure{M^![0]_0|M^{\prime!}[0]_0}
=M(\unsure{0|\omega}).\unsure{L_0^{\Pi_0}|\Pi_0}$.}}}

\margin{This remark about a
$\mathcal{G}^{\unsure{\natural|\sharp}}_0$
now becomes useless.}

These two \inserted{latter} examples suggest the advisability,
in the \inserted{general} case where we would have a subgroup
\deleted{\unsure{$Z$}} $\mathcal{G}^{\unsure{\natural|\sharp}}_0$
of $\mathcal{G}^{\unsure{\natural|\sharp}}$,
\inserted{invariant in $\mathcal{G}$}, such that we have
\[
(16)\qquad
\mathcal{G}^{\unsure{\natural|\sharp}}
=\mathcal{G}^{\unsure{\natural|\sharp}}_0\times\mu
\quad\deleted{\unsure{\subset|\Leftarrow}}
\]
\inserted{(here $\mu$ is the universal central subgroup
$\mu=\{1,\rho^3\}=\{1,\sigma^2\}$)}
of defining $Z_{\rho,\sigma}$ in terms of
$\mathcal{G}^{\unsure{\natural|\sharp}}$, and of
\uline{choosing} the $\alpha,\beta$ of definition (4) in
\deleted{$\unsure{\mathcal{G}^{\natural}_{0\tau}}$}
\[
(17)\qquad
\mathcal{G}^{\unsure{\natural|\sharp}}_{0\tau}
\overset{\mathrm{dfn}}{=}
\mu_\tau.\mathcal{G}^{\unsure{\natural|\sharp}}_0.
\]

\uline{Remark} We are tempted to replace definition (4)
\inserted{(for a) and b))}
by the apparently simpler definition of page 578,
which makes it evident that we have a subgroup.
But they \deleted{\unsure{are compatible}}
\inserted{do not give, even if
$\mathcal{G}^{\unsure{\natural|\sharp}}\subset\Pi$
(\uline{universal case})}
\deleted{with \unsure{if we suppose}}
a trivial action of
\deleted{$\mathcal{G}^{+}/\mathcal{G}^{\unsure{\natural|\sharp}}
\longrightarrow\mathcal{G}_3$}
\inserted{$Z$ on $\mathcal{G}^{+}/\Pi$,
as we \deleted{\unsure{see}}
\inserted{see} immediately,
\deleted{\unsure{$\rho,\sigma$}}}

\mspage{656}

\scanpage{658}

except \unsure{by} \unsure{restricting} ourselves to some, $(u,\mu,\varepsilon_2,\varepsilon_3)$ with $\varepsilon_3=1$,
since $\dot{\rho}^{-1}$ is not equal to $\dot{\rho}$ ! It is true \deleted{that}
(in the universal case for example) that $u(\rho)$ is conjugate \uline{in $\mathcal{G}^{+}$}
to $\rho^{\varepsilon_3}$, but in general it will not be so by an element (say) of $\Pi$.

We set
\[
\begin{gathered}(18)\qquad
L_0^{\unsure{\natural|\sharp}}
=\deleted{\unsure{L_0}\cap\mathcal{G}^{\unsure{\natural|\sharp}}}
\inserted{\bigl\{l\in L_0^{\mathcal{G}}\mid
(\operatorname{int}(l),1,1,1)\in Z_{\rho,\sigma}\bigr\}}
\\[3pt]
\left(\text{where }L_0^{\mathcal{G}}\subset\mathcal{G}^{+}
\text{ is the \unsure{set} of the }\varepsilon_0^n,\ n\in\widehat{\mathbb{Z}}\right)
\end{gathered}\]
\margin{\uline{N.B.} We have
$L_0^{\unsure{\natural|\sharp}}\supset L_0\cap\mathcal{G}^{\unsure{\natural|\sharp}}$,
and the inclusion can be strict, cf. example c) above\ldots
where \uline{$L_0^{\unsure{\natural|\sharp}}=L_0^\Pi$},
$L_0\cap\mathcal{G}^{\unsure{\natural|\sharp}}=L_0^{\Pi_0}$.
\unsure{which} \unsure{is surprising}.}

and we have a canonical homomorphism
\diagram{
\begin{tikzcd}[ampersand replacement=\&]
(19)\quad L_0^{\unsure{\natural|\sharp}}
  \arrow[r,"\unsure{i}"] \& Z_{\rho,\sigma},\\
l \arrow[r,mapsto] \& (\operatorname{int}_{\mathcal{G}^{+}}(l),1,1,1)
\end{tikzcd}
}
injective iff we have
\[
(20)\qquad
L_0^{\unsure{\natural|\sharp}}\cap\operatorname{Centr}(\mathcal{G}^{+})=\{1\}
\]
the image invariant. We set
\[
(21)\qquad
\gamma_{\rho,\sigma}=Z_{\rho,\sigma}/L_0^{\unsure{\natural|\sharp}},
\]
whence
\diagram{
\begin{tikzcd}[ampersand replacement=\&]
(22)\quad 1 \arrow[r] \&
L_0^{\unsure{\natural|\sharp}} \arrow[r] \&
Z_{\rho,\sigma} \arrow[r,"\delta_Z"] \&
\gamma_{\rho,\sigma} \arrow[r] \& 1,
\end{tikzcd}
}
\diagram{
\begin{tikzcd}[ampersand replacement=\&]
 \& \& \mu\\
(23)\quad Z_{\rho,\sigma} \arrow[r] \&
\gamma_{\rho,\sigma}
  \arrow[ur,"\varepsilon_{3\gamma}"]
  \arrow[r,"\chi_\gamma"]
  \arrow[dr,"\varepsilon_{2\gamma}"'] \&
\widehat{\mathbb{Z}}^{*}\\
 \& \& \mu
\end{tikzcd}
}
\unsure{)} \unsure{deduced} from the
\diagram{
\begin{tikzcd}[ampersand replacement=\&, row sep=0pt]
 \& \varepsilon_3\\
(u,\mu,\varepsilon_2,\varepsilon_3)
  \arrow[ur,mapsto] \arrow[r,mapsto] \arrow[dr,mapsto]
  \& \mu\\
 \& \varepsilon_2.
\end{tikzcd}
}

\deleted{
\uline{Remarks} We \deleted{denote} set,
\[
\begin{gathered}(24)\qquad
\left\{
\begin{array}{lll}
\gamma'_{\rho,\sigma}=\operatorname{Ker}\varepsilon_{3\gamma},&
\gamma''_{\rho,\sigma}=\operatorname{Ker}\unsure{\varepsilon_{2,\sigma}|\varepsilon_{2\gamma}},&
\gamma'''_{\rho,\sigma}=\operatorname{Ker}\varepsilon_{2\gamma}\varepsilon_{3\gamma}\\
Z'_{\rho,\sigma}=\ldots&
Z''_{\rho,\sigma}=\ldots&
Z'''_{\rho,\sigma}=\ldots
\end{array}
\right.\\[3pt]
(\text{inverse images in }Z_{\rho,\sigma}\unsure{)}
\end{gathered}\]
}
\mspage{657}

\scanpage{659}
\mspage{632}

We set
\[
(24)\qquad
\left\{
\begin{array}{l}
\gamma'_{\rho,\sigma}=\operatorname{Ker}\varepsilon_{3\gamma},\\[2pt]
\gamma''_{\rho,\sigma}=\operatorname{Ker}\varepsilon_{2\gamma},\\[2pt]
\gamma'''_{\rho,\sigma}=\operatorname{Ker}\varepsilon_{2\gamma}\varepsilon_{3\gamma},\
\gamma^\circ_{\rho,\sigma}=\gamma'\cap\gamma''=\gamma''\cap\gamma'''
=\gamma'''\cap\gamma'\\[4pt]
Z'_{\rho,\sigma},Z''_{\rho,\sigma},Z'''_{\rho,\sigma},Z^\circ_{\rho,\sigma}
\text{ are the inverse images of the }\deleted{\unsure{Z_\rho}}
\gamma'_{\rho,\sigma}\text{ etc}\\
\hfill\text{in }Z_{\rho,\sigma}
\end{array}
\right.
\]
These are thus subgroups of index $1$, \unsure{or} $2$ or four in
$\gamma_{\rho,\sigma}$ \unsure{resp.} $Z_{\rho,\sigma}$ --- the index $4$ being
possible only for $\gamma^\circ_{\rho,\sigma},Z^\circ_{\rho,\sigma}$.

\uline{Remarks} 1) We thus have
\[
(25)\qquad
Z^\circ_{\rho,\sigma}\simeq
\left\{
(u,\mu)\in\operatorname{Aut}(\mathcal{G}^{+})\times\widehat{\mathbb{Z}}^{*}
\left|
\begin{array}{ll}
\mathrm{a})&u(\rho)\text{ is }\mathcal{G}^{\unsure{\natural|\sharp}}\text{-conjugate to }\rho\\
\mathrm{b})&u(\sigma)\quad\uline{\hspace{4em}}\quad\sigma\\
\mathrm{c})&u(\varepsilon_0)=\varepsilon_0^\mu
\end{array}
\right.
\right\}\unsure{,}
\]
\margin{\unsure{More} generally,
\[
Z'''_{\rho,\sigma}\simeq
\left\{
(u,\mu,\varepsilon)\in\operatorname{Aut}(\mathcal{G}^{+})
\times\widehat{\mathbb{Z}}^{*}\times\mu
\left|
\begin{array}{ll}
\mathrm{a})&
u(\rho)\ \deleted{\unsure{\text{is }\varepsilon[\tau](\rho)}}
\inserted{\mathcal{G}^{\unsure{\natural|\sharp}}\text{-conjugate}}
\text{ to }\varepsilon[\tau](\rho)\\
\mathrm{b})&
u(\sigma)\text{ is }\mathcal{G}^{\unsure{\natural|\sharp}}\text{-conjugate to }
\varepsilon[\tau](\sigma)\\
\mathrm{c})&u(\varepsilon_0)=\varepsilon_0^\mu
\end{array}
\right.
\right\}\unsure{,}
\]}
(isomorphism of \deleted{sub}groups). The condition d) of (4) ---
stability of $\mathcal{G}^{\unsure{\natural|\sharp}}$ under $u$ --- is a consequence
here of a) b) c), for we \unsure{will have}
\[
\begin{array}{ll}
u(\rho)\equiv\rho&(\mathcal{G}^{\unsure{\natural|\sharp}})\\
u(\sigma)\equiv\sigma&(\mathcal{G}^{\unsure{\natural|\sharp}})
\end{array}
\]
which implies
\[
(25\mathrm{bis})\qquad
u(x)\equiv x\quad(\mathcal{G}^{\unsure{\natural|\sharp}})
\qquad \forall x\in\mathcal{G}^{+},
\]
i.e. the stability of $\mathcal{G}^{\unsure{\natural|\sharp}}$ \uline{and} a trivial action of
$u$ on $\mathcal{G}/\mathcal{G}^{\unsure{\natural|\sharp}}$.

2) We also have the remarkable element
\[
(26)\qquad
\tau_Z\in Z_{\rho,\sigma},\qquad \tau_Z=(\tau,-1,-1,-1)
\]
\margin{That is one more reason not to require \unsure{in} (4) d) that $u$ operate
trivially on $\mathcal{G}^{+}/\mathcal{G}^{\unsure{\natural|\sharp}}$,
\unsure{which} \unsure{would make me|we would make} lose $\tau$ !}

$3^\circ)$\quad $S_0G_{\rho,\sigma},G_{\rho,\sigma},
S_0G_{\rho,\sigma}^{\unsure{\natural|\sharp}},G_{\rho,\sigma}^{\unsure{\natural|\sharp}}$.

Proceeding as on page 577, we find
\diagram{
\begin{tikzcd}[ampersand replacement=\&]
Z_{\rho,\sigma}\arrow[r]\&\operatorname{Aut}(\mathcal{G}^{+})
\end{tikzcd}
}
hence $1/2$ direct products
\[
(27)\qquad
S_0G_{\rho,\sigma}\overset{\mathrm{dfn}}=
Z_{\rho,\sigma}\mathbin{.}\mathcal{G}^{+},\supset
S_0G_{\rho,\sigma}^{\unsure{\natural|\sharp}}\overset{\mathrm{dfn}}=
Z_{\rho,\sigma}\mathbin{.}\mathcal{G}^{\unsure{\natural|\sharp}}
\]
and a homomorphism
\diagram{
\begin{tikzcd}[ampersand replacement=\&]
(28)\quad L_0^{\unsure{\natural|\sharp}}
  \arrow[r,hook,"\unsure{i_{L_0,S_0G}|i_{L_0,G}}"]\&
S_0G_{\rho,\sigma}^{\unsure{\natural|\sharp}}
  \arrow[r,phantom,"\subset"]\&S_0G_{\rho,\sigma}\\
l\arrow[r,mapsto]\&
\deleted{\unsure{\iota(l)|i(l)}}\,i_{L_0,Z}(l)\mathbin{.}l^{-1}
\end{tikzcd}
}
\mspage{658}

\scanpage{660}

with invariant image, which allows us to set
\[
(29)\qquad
G_{\rho,\sigma}\overset{\mathrm{dfn}}=
\underbrace{S_0G_{\rho,\sigma}}_{Z_{\rho,\sigma}\mathbin{.}\mathcal{G}^{+}}
/\operatorname{Im}L_0^{\unsure{\natural|\sharp}}
\ ,\supset\
G_{\rho,\sigma}^{\unsure{\natural|\sharp}}
=S_0G_{\rho,\sigma}^{\unsure{\natural|\sharp}}
/\operatorname{Im}\unsure{L_0^{\natural}|L_0^{\sharp}}
\]
hence exact sequences
\diagram{
\begin{tikzcd}[ampersand replacement=\&]
(30)\quad 1\arrow[r,dash]\&
L_0^{\unsure{\natural|\sharp}}\arrow[r]\arrow[d,equal]\&
S_0G_{\rho,\sigma}^{\unsure{\natural|\sharp}}\arrow[r]\arrow[d,hook]\&
G_{\rho,\sigma}^{\unsure{\natural|\sharp}}\arrow[r,dash]\arrow[d,hook]\&1,\\
1\arrow[r,dash]\&
L_0^{\unsure{\natural|\sharp}}\arrow[r,dash]\&
S_0G_{\rho,\sigma}\arrow[r]\&G_{\rho,\sigma}\arrow[r,dash]\&1
\end{tikzcd}
}
and a commutative diagram
\diagram{
\begin{tikzcd}[ampersand replacement=\&, column sep=small,row sep=small,nodes in empty cells]
\&Z_{\rho,\sigma}
 \arrow[dr,hook,"i_Z"]
 \arrow[ddr,"{\deleted{\unsure{\longrightarrow}}}" description]
\&\&\\
(31)\quad L_0^{\unsure{\natural|\sharp}}
 \arrow[ur,hook,"i_{L_0,Z}"]
 \arrow[ddr,hook,"i_{L_0,\mathcal{G}}"']
\&\&G_{\rho,\sigma}^{\unsure{\natural|\sharp}}
 \arrow[dr,hook]\&\\[12pt]
\&\&\deleted{\unsure{G_{\rho,\sigma}}}\&G_{\rho,\sigma}\\
\&\mathcal{G}^{\unsure{\natural|\sharp}}
 \arrow[uur,hook,"i_{\mathcal{G}^{\unsure{\natural|\sharp}}}"]
 \arrow[ur,"{\deleted{\unsure{\longrightarrow}}}" description]
 \arrow[dr,hook]\&\&\\
\&\&\mathcal{G}^{+}\arrow[uur,hook,"i_{\mathcal{G}}"']\&
\end{tikzcd}
}
fitting into a commutative diagram of exact sequences
\diagram{
\begin{tikzcd}[ampersand replacement=\&]
(32)\quad 1\arrow[r,dash]\&
L_0^{\unsure{\natural|\sharp}}\arrow[r,dash]\arrow[d,hook]\&
Z_{\rho,\sigma}\arrow[r,"\delta_Z"]\arrow[d,hook]\&
\gamma_{\rho,\sigma}\arrow[r]\arrow[d,Rightarrow]\&1\\
1\arrow[r,dash]\&
\mathcal{G}^{\unsure{\natural|\sharp}}\arrow[r,dash]\arrow[d,hook]\&
G_{\rho,\sigma}^{\unsure{\natural|\sharp}}\arrow[r,"\delta_G"]\arrow[d,hook]\&
\gamma_{\rho,\sigma}\arrow[r]\arrow[d,Rightarrow]\&1.\\
1\arrow[r,dash]\&\mathcal{G}^{+}\arrow[r,dash]\&
G_{\rho,\sigma}\arrow[r]\&\gamma_{\rho,\sigma}\arrow[r]\&1
\end{tikzcd}
}
We consider the composites
\diagram{
\begin{tikzcd}[ampersand replacement=\&]
\&\&\mu\\
(33)\quad G_{\rho,\sigma}\arrow[r,"\delta_G"]\&
\gamma_{\rho,\sigma}
 \arrow[ur,"\varepsilon_{2\gamma}"]
 \arrow[r,"\chi_\gamma"]
 \arrow[dr,"\varepsilon_{3\gamma}"']\&
\widehat{\mathbb{Z}}^{*}\\
\&\&\mu
\end{tikzcd}
}
denoted
\[
\deleted{\unsure{\varepsilon_G=\varepsilon_\gamma\delta_G}}
\]
\[
\chi_G=\chi_\gamma\circ\delta_G,\qquad
\varepsilon_{2G}=\deleted{\unsure{\chi}}\inserted{\varepsilon}_{2\gamma}\delta_G,\qquad
\varepsilon_{3G}=\delta_G\varepsilon_{3\gamma}.
\]

We now have
\[
(34)\qquad
\left\{
\begin{array}{l}
G'_{\rho,\sigma},G''_{\rho,\sigma},G'''_{\rho,\sigma},G^\circ_{\rho,\sigma}
\quad\text{as in }(24)\\
G^+_{\rho,\sigma}=\operatorname{Ker}\chi_G,\quad
G^{'+}_{\rho,\sigma},G^{''+}_{\rho,\sigma},G^{'''+}_{\rho,\sigma},G^{\circ+}_{\rho,\sigma}\\
(\gamma^+_{\rho,\sigma}=\operatorname{Ker}\chi_\gamma\quad
\text{\unsure{denoted}}\ \ldots\unsure{\ldots}\quad)
\end{array}
\right.
\]
\margin{idem \unsure{with} $G_{\rho,\sigma}^{'\unsure{\natural|\sharp}}\ldots$}
If
\[
(35)\qquad
\tau_G\in G_{\rho,\sigma}^{\unsure{\natural|\sharp}}
\inserted{\subset G_{\rho,\sigma}},\quad
\tau_\gamma\in\gamma_{\rho,\sigma}
\]
\mspage{659}

\scanpage{661}
\mspage{633}
denote the images of $\tau_Z$ (26) \unsure{in|in}
$G_{\rho,\sigma},\gamma_{\rho,\sigma}$,
we can make the identification
\[
(36)\qquad
\left\{
\begin{array}{l}
\mathcal G\simeq\text{inverse image of }\{1,\tau_\gamma\}
\inserted{\subset\gamma_{\rho,\sigma}}\deleted{\text{in}}\\
\qquad G_{\rho,\sigma},\quad\text{under }\unsure{G_{\rho,\sigma}}
\longrightarrow\gamma_{\rho,\sigma},\\
\mathcal G^{\unsure{\natural|\sharp}}
=\text{inverse image of }\{1,\tau_\gamma\}\subset\gamma_{\rho,\sigma}
\text{ in }G_{\rho,\sigma}^{\unsure{\natural|\sharp}},
\end{array}
\right.
\]
and
\[
(37)\qquad
\varepsilon_{\mathcal G}
=\chi_G|_{\mathcal G}
=\varepsilon_{2G}|_{\mathcal G}
=\varepsilon_{3G}|_{\mathcal G}
:\left\{
\begin{array}{l}
\text{trivial on }\mathcal G^+\\
-1\text{ on }\tau.
\end{array}
\right.
\]

Let us consider the canonical map
\[
(38)\qquad
\diagram{
\begin{tikzcd}[ampersand replacement=\&, column sep=small,row sep=large]
G_{\rho,\sigma}^{\unsure{\natural|\sharp}}
\arrow[r,hook] \&
G_{\rho,\sigma}^{\deleted{\unsure{\natural|\sharp}}}
\arrow[r] \&
\operatorname{Aut}(\mathcal G^+)\times\widehat{\mathbb Z}^{*}\times\mu\times\mu\\
\&\uline{u}\arrow[r]\&
\bigl(\operatorname{int}(\uline{u})|_{\mathcal G^+},
\chi_G(u),\varepsilon_{2G}(u),\varepsilon_{3G}(u)\bigr)
\end{tikzcd}
}
\]

\uline{Proposition}\quad This is a homomorphism of groups:
\deleted{\unsure{$G_{\rho,\sigma}$}}\deleted{whose}
the image \inserted{of $G_{\rho,\sigma}^{\unsure{\natural|\sharp}}$}
consists of the
\deleted{\unsure{pairs}}\inserted{\unsure{quadruplets|quadruples}}
\[
(u,\mu,\varepsilon_2,\varepsilon_3)\in
\operatorname{Aut}(\mathcal G^+)\times\widehat{\mathbb Z}^{*}\times\mu\times\mu
\]
satisfying the conditions
\[
(39)\qquad
\left\{
\begin{array}{ll}
\text{a)}&
u(\rho)\text{ is }\mathcal G_\tau^{\unsure{\natural|\sharp}}
\text{-conjugate to }\varepsilon_3[\tau]\rho\\
\text{b)}&
u(\sigma)\quad\rule{4.2cm}{0.35pt}\quad
\unsure{\varepsilon_2|\varepsilon_\tau}[\tau]\sigma\\
\text{c)}&
u(\varepsilon_0)\text{ is }\mathcal G^{\unsure{\natural|\sharp}}
\text{-conjugate to }\varepsilon_0^\mu\\
\text{d)}&
u\text{ stabilizes }\mathcal G^{\unsure{\natural|\sharp}}
\quad(\text{automatic in view of a), b) if }\varepsilon_2=\varepsilon_3)
\end{array}
\right.
\]
Its kernel is \deleted{isomorphic to the}
\deleted{\inserted{made up of the elements}}
\[
(40)\quad
\deleted{\unsure{f^{-1}\uline{u}}}
\left\{
f^{-1}\uline{u}\in G_{\rho,\sigma}^{\unsure{\natural|\sharp}}
\ \middle|\
\begin{array}{l}
\deleted{\unsure{f}}\uline{u}=(u,1,1,1)\in Z_{\rho,\sigma}^{\circ+}\\
f\in\mathcal G^{\unsure{\natural|\sharp}}
\end{array}
\quad\left\{\vphantom{\begin{array}{l}u\\u\end{array}}
u=\operatorname{int}(f).
\right.
\right\}.
\]

\deleted{
\uline{Corollary}\quad Let
\[
\unsure{E|\mathcal E}\qquad
(41)\quad
\overset{\inserted{\unsure{\operatorname{Ker}}}}{\unsure{Z_0}}
=\{f\in\mathcal G^{\unsure{\natural|\sharp}}\mid
f(\rho)\text{ is }\mathcal G_\tau^{\unsure{\natural|\sharp}}
\text{-conjugate to }\rho\unsure{\}}.
\]
\[
Z_0^{\unsure{\natural|\sharp}}
=\operatorname{Centr}_{\unsure{\mathcal G^{\natural}|\mathcal G}}
(\varepsilon_0\unsure{)}.
\]
}

\uline{Corollary}\quad Let
\[
(41)\qquad
Z_0=\operatorname{Centr}_{\mathcal G^+}(L_0),\qquad
Z_0^{\unsure{\natural|\sharp}}
=Z_0\cap\mathcal G^{\unsure{\natural|\sharp}}
\]
\[
\inserted{
\left(=\{f\in\mathcal G^{\unsure{\natural|\sharp}}\mid
(\operatorname{int}(f),1,1,1)\in
\unsure{Z_{\rho,\sigma}^{\circ+}|Z_{\rho,\sigma}}\unsure{\}}\unsure{)}\right.
}
\]
so that $L_0^{\unsure{\natural|\sharp}}$ is identified with a subgroup
\mspage{660}

\scanpage{662}
central in $Z_0^{\unsure{\natural|\sharp}}$.
Let us consider the map
\[
(42)\qquad
\diagram{
\begin{tikzcd}[ampersand replacement=\&, column sep=small,row sep=large,column sep=9pt,row sep=14pt,execute at begin picture={\fontsize{10}{12}\selectfont}]
Z_0^{\unsure{\natural|\sharp}}\arrow[rr]\&\&
S_0G_{\rho,\sigma}^{\unsure{\natural|\sharp}}
=Z_{\rho,\sigma}\cdot\mathcal G^{\unsure{\natural|\sharp}}
\&\&\\
f\arrow[r,mapsto]\&
f^{-1}\arrow[d,phantom,"\in"{rotate=-90}]\&
{\cdot\uline{u}(f)}\arrow[d,phantom,"\in"{rotate=-90}]\&
\inserted{=\uline{u}(f)\cdot f^{-1}}\&
{\text{where }\uline{u}(f)=(\operatorname{int}(f),1,1,1)
\in Z_{\rho,\sigma}}\arrow[d,phantom,"\subset"{rotate=-90}]\\
\&\mathcal G^{\unsure{\natural|\sharp}}\&
Z_{\rho,\sigma}\&\&G_{\rho,\sigma}^{\unsure{\natural|\sharp}}
\end{tikzcd}
}
\]
This map is a homomorphism of groups,
\deleted{\unsure{whose|definition}}\inserted{it}
\deleted{\inserted{inverse image of
$\unsure{Z_0|L_0}\ \deleted{\unsure{L_0}}$}}
\deleted{\unsure{induces|injective} \unsure{by}
\unsure{projection|passage} \unsure{to the} \unsure{quotient}
$L_0^{\unsure{\natural|\sharp}}$}
\deleted{the image}
induces on $L_0^{\unsure{\natural|\sharp}}\subset
Z_0^{\unsure{\natural|\sharp}}$ the canonical homomorphism
$\inserted{\unsure{i_{L_0,SG}|i_{L_0,S_0G}}}$
(28), and $L_0^{\unsure{\natural|\sharp}}\subset
Z_0^{\unsure{\natural|\sharp}}$ is the inverse image of
\[
\unsure{i_{L_0,SG}|i_{L_0,S_0G}}
(L_0^{\unsure{\natural|\sharp}})
=\operatorname{Ker}
\bigl(S_0G_{\rho,\sigma}^{\unsure{\natural|\sharp}}
\longrightarrow G_{\rho,\sigma}^{\unsure{\natural|\sharp}}\bigr).
\]
We thus find an injective homomorphism
\[
(43)\qquad
\diagram{
\begin{tikzcd}[ampersand replacement=\&]
Z_0^{\unsure{\natural|\sharp}}/L_0^{\unsure{\natural|\sharp}}
\arrow[r,hook]\&G_{\rho,\sigma}^{\unsure{\natural|\sharp}}
\end{tikzcd}
}
\]
whose image is the kernel of (38), i.e. we have a canonical exact sequence
\[
(44)\qquad
\diagram{
\begin{tikzcd}[ampersand replacement=\&, column sep=small,row sep=large]
1\arrow[r]\&
L_0^{\unsure{\natural|\sharp}}\arrow[r]\&
Z_0^{\unsure{\natural|\sharp}}\arrow[r]\&
G_{\rho,\sigma}^{\unsure{\natural|\sharp}}\arrow[r]\&
\widetilde G_{\rho,\sigma}^{\unsure{\natural|\sharp}}
\arrow[r]\arrow[d,phantom,"\subset"{rotate=-90}]\&1\\
\&\&\&\&
{\operatorname{Aut}(\mathcal G^+)\times\widehat{\mathbb Z}^{*}\times\mu\times\mu}
\&
\end{tikzcd}
}
\]
where $\widetilde G_{\rho,\sigma}^{\unsure{\natural|\sharp}}$
consists of the $(u,\mu,\varepsilon_2,\varepsilon_3)$ satisfying
the conditions (39).

\uline{Remark.}\quad This proposition brings out
to what extent the \deleted{\illegible{1}} construction made here of
$G_{\rho,\sigma}^{\unsure{\natural|\sharp}}$
is finer than the one \unsure{where}
we \unsure{would define|defined} it simply as a
group of automorphisms of $\mathcal G^+$
(we even \unsure{have|of them} need
of the $\mu,\varepsilon,\varepsilon'$
into the \unsure{bargain}\ldots).
\mspage{661}
